\section{控制边界框}
引自 \tkzimp{PgfManual}：

"当添加 clip 选项时，当前路径将用于裁剪后续的图形。裁剪绝不会扩大裁剪区域。因此，当你针对某一路径裁剪后再针对另一路径裁剪时，实际裁剪区域是两者的交集。
唯一能扩大裁剪路径的方法是结束进行裁剪的 {pgfscope}。在 {pgfscope} 结束时，进入该作用域时生效的裁剪路径会被重新恢复。"


首先，使用 \TIKZ\ 时你无需处理边界框的大小。早期版本的 \tkzNamePack{tkz-euclide} 并不控制它；自第 4 版起，边界框的大小受到了限制。

使用宏 \tkzcname{tkzInit} 后，初始边界框由以 $(0,0)$ 和 $(10,10)$ 为对角顶点的矩形定义。宏 \tkzcname{tkzInit} 可以通过参数 \tkzname{xmin}、\tkzname{xmax}、\tkzname{ymin} 和 \tkzname{ymax} 修改这个初始边界框。当然，任何画在外部的图形都会修改由 \TIKZ\ 维护的边界框。你可以通过 \TIKZ\ 命令或选项来影响这一行为，例如 \tkzcname{useasboundingbox} 或选项 \tkzname{use as bounding box}。这样你可以为图形预留一个框，但图形可能溢出该框并延伸到正文中。
命令 \tkzcname{pgfresetboundingbox} 会清除边界框并重新开始一个新的边界框。

\subsection{\tkzcname{tkzInit} 的作用}
 然而，有时需要控制所显示内容的大小。
 为此，你需要先准备好将要使用的边界框：这正是宏 \tkzNameMacro{tkzInit} 的作用。对于某些图形，固定极值（xmin、xmax、ymin 和 ymax）并对该矩形执行 \code{clip}，以便尽可能精确地控制图形大小，是很有用的。

有两个宏可用于控制边界框：
\begin{itemize}
   \item \tkzcname{tkzInit}
   \item \tkzcname{tkzClip}
\end{itemize}
\vspace{20pt}

我还添加了一些与边界框直接相关的宏：你可以显示它或对其进行裁剪（见下文）。

\subsection{\tkzcname{tkzInit}}

\begin{NewMacroBox}{tkzInit}{\oarg{local options}}\hypertarget{init}{}%
\begin{tabular}{lll}%    
选项  & 默认值 & 说明             \\
\midrule    
\TOline{xmin} {0} {横坐标的最小值，单位为 cm}
\TOline{xmax} {10} {横坐标的最大值，单位为 cm}
\TOline{xstep}{1} {$x$ 方向两个刻度之间的间距}
\TOline{ymin} {0} {$y$ 轴的最小值，单位为 cm}
\TOline{ymax} {10} {$y$ 轴的最大值，单位为 cm}
\TOline{ystep}{1} {$y$ 方向两个刻度之间的间距}  
\bottomrule    
\end{tabular}

\medskip 

\tkzcname{tkzInit} 的作用是定义一个 \textcolor{red}{正交}坐标系和平面上的一个矩形区域，你可以在该区域内用笛卡尔坐标放置图形。
该宏定义了你的工作区域，就像在图形计算器上一样。自 \tkzname{\tkznameofpack} 第 4 版起，\tkzcname{xstep} 和 \tkzcname{ystep} 始终为 1，因此 \tkzcname{tkzInit} 主要用于 \code{clipping out} 之类的操作。
\end{NewMacroBox}


\subsection{\tkzcname{tkzClip}}

\begin{NewMacroBox}{tkzClip}{\oarg{local options}}
该宏会隐藏位于由 (xmin~;~ymin) 与 (xmax~;~ymax) 所定义矩形之外的所有内容。

\medskip
\begin{tabular}{lll}
\hline
选项  & 默认值 & 说明             \\
\midrule
\TOline{space} {1} {在背景矩形的右、左、下、上四个方向上增加的数值}
\bottomrule
\end{tabular}

\medskip

选项 \tkzname{space} 会扩大图形的可见部分，使其变为由 (xmin-space~;~ymin-space) 与 (xmax+space~;~ymax+space) 所定义的矩形。\tkzname{space} 可以为负值！单位为 cm，且无需写出。
\end{NewMacroBox}



该宏对初始矩形进行 \code{clip}，使得只有位于该矩形内部的路径才会被绘制。

\begin{tkzexample}[latex=8cm,small]
\begin{tikzpicture}
 \tkzInit[xmax=4, ymax=3]
 \tkzDefPoints{-1/-1/A,5/2/B}
 \tkzDrawX \tkzDrawY 
 \tkzGrid
 \tkzClip
 \tkzDrawSegment(A,B)
\end{tikzpicture}
\end{tkzexample} 

你可以增加一些空间：
\begin{tkzltxexample}[]
  \tkzClip[space=1]
\end{tkzltxexample} 

\subsection{\tkzcname{tkzClip} 与选项 \tkzname{space}} 
该选项在 \code{clipped} 矩形周围增加一些空间。
\begin{tkzexample}[latex=8cm,small]
\begin{tikzpicture}
 \tkzInit[xmax=4, ymax=3]
 \tkzDefPoints{-1/-1/A,5/2/B}
 \tkzDrawX \tkzDrawY 
 \tkzGrid
 \tkzClip[space=1]
 \tkzDrawSegment(A,B)
\end{tikzpicture}
\end{tkzexample}   
\code{clipped} 矩形的尺寸为 \tkzname{xmin-1}、\tkzname{ymin-1}、\tkzname{xmax+1} 和 \tkzname{ymax+1}。 

%<--------------------------------------------------------------------------->
%              tkzShowBB
%<--------------------------------------------------------------------------->
\subsection{\tkzcname{tkzShowBB}}
最简单的宏。
\begin{NewMacroBox}{tkzShowBB}{\oarg{local options}}% 
该宏以矩形框的形式显示边界框。它接受 \TIKZ\ 选项。
\end{NewMacroBox} 


\subsubsection{\tkzcname{tkzShowBB} 示例}
\begin{tkzexample}[latex=8cm,small]
\begin{tikzpicture}[scale=.5]
  \tkzInit[ymax=5,xmax=8]
  \tkzGrid  
  \tkzDefPoint(3,0){A}
   \begin{scope}
    \tkzClipBB
    \tkzDefCircle[R](A,5) \tkzGetPoint{a}
    \tkzDrawCircle(A,a)
    \tkzShowBB[line width = 4pt,fill=teal!10,%
              opacity=.4]
   \end{scope}
\tkzDefCircle[R](A,4) \tkzGetPoint{b}
\tkzDrawCircle[red](A,b)
\end{tikzpicture}
\end{tkzexample}
%<--------------------------------------------------------------------------->
%         tkzClipBB
%<--------------------------------------------------------------------------->
\subsection{\tkzcname{tkzClipBB}}
\begin{NewMacroBox}{tkzClipBB}{}%
其思想是将后续的构造限制在当前边界框之内。
\end{NewMacroBox}

\subsubsection{\tkzcname{tkzClipBB} 与角平分线示例}

\begin{tkzexample}[latex=6cm,small]
  \begin{tikzpicture}
  \tkzInit[xmin=-3,xmax=6, ymin=-1,ymax=6]
  \tkzDefPoint(0,0){O}\tkzDefPoint(3,1){I}
  \tkzDefPoint(1,4){J}
  \tkzDefLine[bisector](I,O,J) \tkzGetPoint{i}
  \tkzDefLine[bisector out](I,O,J) \tkzGetPoint{j}
  \tkzDrawPoints(O,I,J,i,j)
  \tkzClipBB
  \tkzDrawLines[add = 1 and 2,color=orange](O,I O,J)
  \tkzDrawLines[add = 1 and 2](O,i O,j)
  \tkzShowBB
  \end{tikzpicture}
\end{tkzexample}


\newpage

\section{裁剪不同对象}

\subsection{裁剪多边形} 
 \begin{NewMacroBox}{tkzClipPolygon}{\oarg{local options}\parg{points list}}%
该宏将图形限制在给定的多边形之内。

\medskip
\begin{tabular}{lll}%
\toprule
参数       & 示例 & 说明     \\ 
\midrule
\TAline{\parg{pt1,pt2,pt3,\dots}}{\parg{A,B,C}}{}
\midrule
选项  & 默认值 & 说明             \\
\midrule    
\TOline{out} {} {裁剪对象的外部}
 \end{tabular}
\end{NewMacroBox}

\subsubsection{\tkzcname{tkzClipPolygon}}

\begin{tkzexample}[latex=7cm,small]
  \begin{tikzpicture}[scale=1.25] 
  \tkzDefPoint(0,0){A} 
  \tkzDefPoint(4,0){B} 
  \tkzDefPoint(1,3){C} 
  \tkzDrawPolygon(A,B,C) 
  \tkzDefPoint(0,2){D} 
  \tkzDefPoint(2,0){E} 
  \tkzDrawPoints(D,E) 
  \tkzLabelPoints(D,E) 
  \tkzClipPolygon(A,B,C) 
  \tkzDrawLine[new](D,E)
\end{tikzpicture}
\end{tkzexample}

\subsubsection{\tkzcname{tkzClipPolygon[out]}}

\begin{tkzexample}[latex=7cm,small]
  \begin{tikzpicture}[scale=1]
  \tkzDefPoint(0,0){P1}
  \tkzDefPoint(4,0){P2}
  \tkzDefPoint(4,4){P3}
  \tkzDefPoint(0,4){P4}
  \tkzDefPoint(1,1){Q1}
  \tkzDefPoint(3,1){Q2}
  \tkzDefPoint(3,3){Q3}
  \tkzDefPoint(1,3){Q4}
  \tkzDrawPolygon(P1,P2,P3,P4)
  \begin{scope}
     \tkzClipPolygon[out](Q1,Q2,Q3,Q4)
    \tkzFillPolygon[teal!20](P1,P2,P3,P4)
  \end{scope}
  \tkzDrawPolygon(Q1,Q2,Q3,Q4)
  \end{tikzpicture}
\end{tkzexample}

\subsubsection{示例：在正方形中使用 \code{Clip} 绘制 Sangaku} 
\begin{tkzexample}[latex=7cm, small]  
\begin{tikzpicture}[scale=.75]
 \tkzDefPoint(0,0){A} \tkzDefPoint(8,0){B}
 \tkzDefSquare(A,B)   \tkzGetPoints{C}{D}
 \tkzDefPoint(4,8){F}
 \tkzDefTriangle[equilateral](C,D) 
 \tkzGetPoint{I}
 \tkzDefPointBy[projection=onto B--C](I) 
 \tkzGetPoint{J}
 \tkzInterLL(D,B)(I,J)  \tkzGetPoint{K}
 \tkzDefPointBy[symmetry=center K](B)  
 \tkzGetPoint{M}
 \tkzClipPolygon(B,C,D,A)
 \tkzFillPolygon[color = orange](A,B,C,D)
 \tkzFillCircle[color = yellow](M,I)
 \tkzFillCircle[color = blue!50!black](F,D)
\end{tikzpicture}
\end{tkzexample}

\subsection{裁剪圆盘}

\begin{NewMacroBox}{tkzClipCircle}{\oarg{local options}\parg{A,B}}%
\begin{tabular}{lll}%
\toprule
参数           & 示例 & 说明                         \\
\midrule
\TAline{\parg{A,B}}{\parg{A,B}} {AB 为半径}
\bottomrule
\end{tabular}  
 
\medskip
\begin{tabular}{lll}%
选项             & 默认值 & 说明                         \\ 
\midrule
\TOline{out} {} {裁剪对象的外部}
 \bottomrule
\end{tabular}

\medskip
无需写出 \tkzname{radius}，因为它是默认选项。
\end{NewMacroBox}

 \subsubsection{简单裁剪} 
\begin{tkzexample}[latex=6cm,small] 
\begin{tikzpicture}[scale=.5]
  \tkzDefPoint(0,0){A} \tkzDefPoint(2,2){O}
  \tkzDefPoint(4,4){B} \tkzDefPoint(5,5){C}
  \tkzDrawPoints(O,A,B,C) 
  \tkzLabelPoints(O,A,B,C)
  \tkzDrawCircle(O,A) 
  \tkzClipCircle(O,A)
  \tkzDrawLine(A,C)
  \tkzDrawCircle[fill=teal!10,opacity=.5](C,O)
\end{tikzpicture} 
\end{tkzexample}

\subsection{向外裁剪}

\begin{tkzexample}[latex=6cm,small]
\begin{tikzpicture}
  \tkzInit[xmin=-3,ymin=-2,xmax=4,ymax=3]
   \tkzDefPoint(0,0){O}
   \tkzDefPoint(-4,-2){A}
   \tkzDefPoint(3,1){B}
   \tkzDefCircle[R](O,2) \tkzGetPoint{o}
   \tkzDrawPoints(A,B) % to have a good bounding box
   \begin{scope}
    \tkzClipCircle[out](O,o)
    \tkzDrawLines(A,B)
   \end{scope}
\end{tikzpicture}  
\end{tkzexample} 

\subsection{圆盘的交集}

\begin{tkzexample}[latex=6cm,small]
\begin{tikzpicture}
\tkzDefPoints{0/0/O,4/0/A,0/4/B}
\tkzDrawPolygon[fill=teal](O,A,B)
\tkzClipPolygon(O,A,B)
\tkzClipCircle(A,O)
\tkzClipCircle(B,O)
\tkzFillPolygon[white](O,A,B)
\end{tikzpicture}
\end{tkzexample} 

关于裁剪圆的更复杂示例，请参见 \ref{About clipping circles} 节。

\subsection{裁剪扇形}
\tkzHandBomb\  注意，参数取决于所使用的选项。
\begin{NewMacroBox}{tkzClipSector}{\oarg{local options}\parg{O,\dots}\parg{\dots}}%
\begin{tabular}{lll}%
选项             & 默认值 & 说明                         \\ 
\midrule
\TOline{towards}{towards}{$O$ 为圆心，扇形从 $A$ 起沿 $(OB)$ 方向}
\TOline{rotate} {towards}{扇形从 $A$ 起，由角度确定其张角} 
\TOline{R}{towards}{给定半径和两个角度} 
\bottomrule
\end{tabular} 

\medskip
当然，你可以为图形添加所有 \TIKZ\ 样式。

\medskip   
\begin{tabular}{lll}%
\toprule
选项             & 参数 & 示例                         \\ 
\midrule
\TOline{towards}{\parg{pt,pt}\parg{pt}}{\tkzcname{tkzClipSector(O,A)(B)}}
\TOline{rotate} {\parg{pt,pt}\parg{angle}}{\tkzcname{tkzClipSector[rotate](O,A)(90)}} 
\TOline{R}{\parg{pt,$r$}\parg{angle 1,angle 2}}{\tkzcname{tkzClipSector[R](O,2)(30,90)}}
\end{tabular}
\end{NewMacroBox}

\subsubsection{裁剪扇形：第一个示例}
\begin{tkzexample}[latex=7cm,small] 
\begin{tikzpicture}[scale=0.5]
\tkzDefPoint(0,0){a}
\tkzDefPoint(12,0){b}
\tkzDefPoint(4,10){c}
\tkzInterCC[R](a,6)(b,8) 
\tkzGetFirstPoint{AB1} \tkzGetSecondPoint{AB2}
\tkzInterCC[R](a,6)(c,6) 
\tkzGetFirstPoint{AC1} \tkzGetSecondPoint{AC2}
\tkzInterCC[R](b,8)(c,6) 
\tkzGetFirstPoint{BC1} \tkzGetSecondPoint{BC2}
\tkzDrawArc(a,AB2)(AB1)
\tkzDrawArc(b,AB1)(AB2)
\tkzDrawArc(a,AC2)(AC1)
\tkzDrawArc(c,AC1)(AC2)
\tkzDrawArc(b,BC2)(BC1)
\tkzDrawArc(c,BC1)(BC2)
\begin{scope}
\tkzClipSector(b,BC2)(BC1)
\tkzFillSector[teal!40!white](c,BC1)(BC2)
\end{scope}
\begin{scope}
\tkzClipSector(a,AB2)(AB1)
\tkzFillSector[teal!40!white](b,AB1)(AB2)
\end{scope}
\begin{scope}
\tkzClipSector(a,AC2)(AC1)
\tkzFillSector[teal!40!white](c,AC1)(AC2)
\end{scope}
\end{tikzpicture}
\end{tkzexample}

\subsubsection{裁剪扇形：第二个示例}
\begin{tkzexample}[latex=7cm,small] 
\begin{tikzpicture}[scale=1.5] 
  \tkzDefPoint(0,0){O}
  \tkzDefPoint(2,-1){A}
  \tkzDefPoint(1,1){B} 
  \tkzDrawSector[new,dashed](O,A)(B)
  \tkzDrawSector[new](O,B)(A)
\begin{scope}
\tkzClipSector(O,B)(A)
\tkzDefSquare(O,B) \tkzGetPoints{B'}{O'}
\tkzDrawPolygon[color=teal,fill=teal!20](O,B,B',O')
\end{scope}
\tkzDrawPoints(A,B,O) 
\end{tikzpicture} 
\end{tkzexample}


\subsection{\TIKZ\ 选项：trim left 或 trim right}
参见 \tkzimp{pgfmanual}。

\subsection{\TIKZ\ 控制 \tkzcname{pgfinterruptboundingbox} 与 \tkzcname{endpgfinterruptboundingbox}}
该命令会临时中断边界框的计算并开始一个新的边界框。
参见 \tkzimp{pgfmanual}

\subsubsection{控制边界框} 
\begin{tkzexample}[latex=7cm,small] 
\begin{tikzpicture}
\tkzDefPoint(0,5){A}\tkzDefPoint(5,4){B}
\tkzDefPoint(0,0){C}\tkzDefPoint(5,1){D}
\tkzDrawSegments(A,B C,D A,C)
\pgfinterruptboundingbox
   \tkzInterLL(A,B)(C,D)\tkzGetPoint{I}
\endpgfinterruptboundingbox
\tkzClipBB
\tkzDrawCircle(I,B)
\end{tikzpicture}
\end{tkzexample}

\subsection{反向裁剪：\tkzname{tkzreverseclip}}

要使用该选项，必须先定义一个边界框。

\begin{tkzltxexample}[]
  \tikzset{tkzreverseclip/.style={insert path={
     (current bounding box.south west) --(current bounding box.north west)
   --(current bounding box.north east) --  (current bounding box.south east)
   -- cycle} }}
\end{tkzltxexample}

\subsubsection{\tkzcname{tkzClipPolygon[out]} 示例}
\tkzcname{tkzClipPolygon[out]}、\tkzcname{tkzClipCircle[out]} 使用了该选项。
\begin{tkzexample}[vbox,small]
\begin{tikzpicture}[scale=1]
\tkzInit[xmin=-5,xmax=5,ymin=-4,ymax=6]
\tkzClip
\tkzDefPoints{-.5/0/P1,.5/0/P2}
\foreach \i [count=\j from 3] in {2,...,7}{%
    \tkzDefShiftPoint[P\i]({45*(\i-1)}:1){P\j}}  
\tkzClipPolygon[out](P1,P...,P8)
\tkzCalcLength(P1,P5)\tkzGetLength{r}
\begin{scope}[blend group=screen]
  \foreach \i in {1,...,8}{%
   \tkzDefCircle[R](P\i,\r) \tkzGetPoint{x}
   \tkzFillCircle[color=teal](P\i,x)}
  \end{scope}
\end{tikzpicture}
\end{tkzexample}

\endinput
